\documentclass[11pt,letterpaper]{article}
\def\SuiteDocTitle{Lab 03: Slice a full workflow vertically}
\def\SuiteDocCode{LAB-03}
% Shared professional style: pdfLaTeX compatible, no external font files.
\usepackage[T1]{fontenc}
\usepackage[utf8]{inputenc}
\usepackage{lmodern}
\usepackage{microtype}
\usepackage[margin=0.83in,headheight=16pt,footskip=28pt]{geometry}
\usepackage[table]{xcolor}
\definecolor{navy}{HTML}{12304A}
\definecolor{teal}{HTML}{087F83}
\definecolor{gold}{HTML}{B78633}
\definecolor{blueback}{HTML}{F0F5F8}
\definecolor{tealback}{HTML}{EEF8F6}
\definecolor{amberback}{HTML}{FFF7E8}
\definecolor{grayback}{HTML}{F6F7F8}
\usepackage{array,tabularx,longtable,booktabs}
\usepackage{enumitem}
\usepackage{amssymb}
\usepackage{titlesec}
\usepackage{fancyhdr,lastpage}
\usepackage{tcolorbox}
\tcbuselibrary{breakable}
\usepackage{pgffor}
\usepackage{hyperref}
\hypersetup{colorlinks=true,linkcolor=teal,urlcolor=teal,pdfauthor={Engineering Learning Series},pdftitle={\SuiteDocTitle}}
\setlength{\parindent}{0pt}
\setlength{\parskip}{5.5pt}
\setlist{topsep=3pt}
\setcounter{tocdepth}{2}
\titleformat{\section}{\Large\bfseries\color{navy}}{\thesection}{0.7em}{}
\titleformat{\subsection}{\large\bfseries\color{teal}}{\thesubsection}{0.6em}{}
\titleformat{\subsubsection}{\normalsize\bfseries\color{navy}}{\thesubsubsection}{0.6em}{}
\pagestyle{fancy}
\fancyhf{}
\fancyhead[L]{\scriptsize\color{navy}AGILE REQUIREMENTS ENGINEERING}
\fancyhead[R]{\scriptsize\color{teal}\SuiteDocCode}
\fancyfoot[L]{\scriptsize Independent learning and implementation materials}
\fancyfoot[R]{\scriptsize Page \thepage\ of \pageref{LastPage}}
\renewcommand{\headrulewidth}{0.35pt}
\renewcommand{\footrulewidth}{0pt}
\newcommand{\SuiteHeader}[2]{%
{\color{teal}\small\bfseries ENGINEERING DOCUMENTATION SUITE \hfill #2}\par\vspace{8pt}
{\color{navy}\LARGE\bfseries\SuiteDocTitle}\par\vspace{5pt}
{\color{black!68}\normalsize #1}\par\vspace{8pt}
{\color{gold}\rule{\linewidth}{2pt}}\par\vspace{8pt}}
\newcommand{\WritingLines}[1]{\par\foreach \n in {1,...,#1}{\noindent\rule{\linewidth}{0.35pt}\par\vspace{0.28cm}}}
\newcommand{\FreeSpace}[1]{\begin{tcolorbox}[colback=white,colframe=black!25,height=#1,boxrule=0.5pt]\end{tcolorbox}}
\newcommand{\Field}[1]{\noindent\textbf{#1}\quad\rule{0.42\linewidth}{0.35pt}\par\vspace{0.3cm}}
\newcommand{\CheckItem}[1]{\noindent$\square$\quad #1\par\vspace{2pt}}
\newcommand{\SmallNote}[1]{{\footnotesize\color{black!62}#1}\par}
\newcommand{\Guidance}[1]{\begin{tcolorbox}[colback=blueback,colframe=teal,breakable,title=Guidance]#1\end{tcolorbox}}
\sloppy

\begin{document}
\SuiteHeader{Thin increments and dependency exposure | Practical assessment}{Hands-on laboratory}
\tableofcontents\medskip\noindent\rule{\linewidth}{0.4pt}\medskip
\section{Laboratory brief}
\begin{longtable}{|p{0.28\linewidth}|p{0.60\linewidth}|}\hline
\textbf{Attribute} & \textbf{Value} \\ \hline\endfirsthead\hline
\textbf{Attribute} & \textbf{Value} \\ \hline\endhead
Focus & Thin increments and dependency exposure \\ \hline
Prerequisite reading & SG-02, SG-07, SG-10 \\ \hline
Difficulty & Practitioner \\ \hline
Suggested duration & 60–90 minutes \\ \hline
\end{longtable}
\section{Scenario and initial inputs}
An epic asks for request submission, attachment upload, multi-stage review, reminders, audit logging and analytics in one delivery.

\section{Procedure}
\begin{enumerate}[leftmargin=*,itemsep=3pt]
\item Draw a one-line end-to-end journey from requester action to feedback.
\item Produce four incremental slices each yielding a demonstrable result.
\item Define a vertical integration test for the first slice.
\item Identify dependencies and one technical experiment that reduces major uncertainty.
\item Explain why UI-only and database-only stories are inappropriate slices.
\end{enumerate}
\section{Submission checklist}
\begin{itemize}[leftmargin=*,itemsep=2pt]
\item An epic slicing diagram or table
\item Four value-bearing slices
\item A first-slice demonstration script
\end{itemize}
\section{Verification criteria}
\begin{itemize}[leftmargin=*,itemsep=2pt]
\item Outputs describe user outcomes rather than only technical activity.
\item Evidence is sufficient for another practitioner to reproduce or challenge the work.
\item Security, negative behavior and assumptions are visible when relevant.
\item Explicit tradeoffs and decision rights appear in the results.
\end{itemize}
\begin{tcolorbox}[colback=amberback,colframe=navy,title={Instructor note / expected reasoning},breakable]
A first iteration may have a text-only form and persistence, but must produce a usable confirmation.

\end{tcolorbox}
\section{Stretch exercise}
Exchange artifacts with another participant. Identify one unverifiable expectation, one dependency hidden by the initial plan, and one incomplete traceability link. Revise the record and document what changed.

\section{Evidence capture sheet}
\Field{Participant / team}\Field{Lab date}\Field{Repository / artifact reference}
\WritingLines{7}

\end{document}
